\documentclass[12pt]{article}
\usepackage{amsmath,amssymb}

\begin{document}
\section*{{2021 USAMO Problems/Problem 3}}
Source: AoPS Wiki

\subsection*{Problem}
Let $n \geq 2$ be an integer. An $n \times n$ board is initially empty. Each minute, you may perform one of three moves: If there is an L-shaped tromino region of three cells without stones on the board (see figure; rotations not allowed), you may place a stone in each of those cells. If all cells in a column have a stone, you may remove all stones from that column. If all cells in a row have a stone, you may remove all stones from that row. [asy] unitsize(20); draw((0,0)--(4,0)--(4,4)--(0,4)--(0,0)); fill((0.2,3.8)--(1.8,3.8)--(1.8, 1.8)--(3.8, 1.8)--(3.8, 0.2)--(0.2, 0.2)--cycle, grey); draw((0.2,3.8)--(1.8,3.8)--(1.8, 1.8)--(3.8, 1.8)--(3.8, 0.2)--(0.2, 0.2)--(0.2, 3.8), linewidth(2)); draw((0,2)--(4,2)); draw((2,4)--(2,0)); [/asy] For which $n$ is it possible that, after some non-zero number of moves, the board has no stones? Solution 1 Label the bottom right cell as $\omega^0$ where $\omega = e^{i2\pi/3}$ . Then label each cell with $\omega^d$ where $d$ is the Manhattan distance of the current cell from the bottom right cell. No matter where an L-shaped tromino is placed, the sum of the labels in the cells is, \[\omega^0 + \omega^1 + \omega^2 = \frac{\omega^3 - 1}{\omega - 1} = \frac{0}{\omega - 1} = 0\] If $n = 3k + 1$ . Then the sum of all the labels in the cells in a row or column is, \[k(\omega^0 + \omega^1 + \omega^2)+\omega^0 = 1\] If $n = 3k + 2$ . Then the sum of all the labels in cells in a row or column is, \[k(\omega^0 + \omega^1 + \omega^2) + \omega^0 + \omega^1 = 1 + \omega\] Let $S$ be the sum of all the labels of cells containing a stone. In either above case, placing an L-shaped Tromino leaves $S$ invariant but removing a row or column of stone causes $\textrm{Re}(S)$ to decrease (by $1$ or $1/2$ in each case). However the empty grid has $S=0$ , so this can't be reached. Hence $n=3k$ . We now provide a sequence of moves which empties the grid. Divide the grid into $3\times 3$ grids and do the following to every sub-grid. (When I refer to placing an L-shaped tromino on a cell, I am refering to the position of the bottom left stone) Add an L-shaped tromino to the bottom left cell and middle cell Remove the middle row Add an L-shaped tromino to the middle left cell Remove the left and middle columns and we're done, the grid is empty Let $n \geq 2$ be an integer. An $n \times n$ board is initially empty. Each minute, you may perform one of three moves: If there is an L-shaped tromino region of three cells without stones on the board (see figure; rotations not allowed), you may place a stone in each of those cells. If all cells in a column have a stone, you may remove all stones from that column. If all cells in a row have a stone, you may remove all stones from that row. [asy] unitsize(20); draw((0,0)--(4,0)--(4,4)--(0,4)--(0,0)); fill((0.2,3.8)--(1.8,3.8)--(1.8, 1.8)--(3.8, 1.8)--(3.8, 0.2)--(0.2, 0.2)--cycle, grey); draw((0.2,3.8)--(1.8,3.8)--(1.8, 1.8)--(3.8, 1.8)--(3.8, 0.2)--(0.2, 0.2)--(0.2, 3.8), linewidth(2)); draw((0,2)--(4,2)); draw((2,4)--(2,0)); [/asy] For which $n$ is it possible that, after some non-zero number of moves, the board has no stones?

\end{document}
